\section{Introduction}

\label{sec:intro}
% ============================================================================

Cryptographically relevant quantum computers (CRQCs) threaten the security
foundations of all modern infrastructure. Shor's algorithm~\cite{shor1994}
breaks RSA, ECDSA, and ECDH in polynomial time. Grover's
algorithm~\cite{grover1996} weakens symmetric ciphers and hash functions by a
quadratic factor.

Hanzo infrastructure relies on cryptography at every layer:
\begin{itemize}[leftmargin=1.1em]
    \item \textbf{TLS}: ECDHE key exchange (X25519) and ECDSA server
    certificates for all service-to-service and client-to-service connections.
    \item \textbf{IAM}: Ed25519 JWT signing for authentication tokens.
    \item \textbf{KMS}: AES-256-GCM envelope encryption with ECDH-derived
    key wrapping for secret storage.
    \item \textbf{Agent credentials}: Ed25519 certificate chains for delegated
    agent permissions.
    \item \textbf{ACI}: ECDSA (secp256k1) for blockchain transaction signing
    and compute attestations.
\end{itemize}

While large-scale CRQCs do not yet exist, the ``harvest now, decrypt later''
threat model requires action today: adversaries recording encrypted traffic can
decrypt it once quantum computers become available. Data with long
confidentiality requirements (medical records, financial data, government
communications) processed through Hanzo AI must be protected against future
quantum decryption.

\subsection{NIST Standardization}

NIST finalized three post-quantum standards in August 2024~\cite{nistpq}:
\begin{enumerate}[leftmargin=1.4em]
    \item \textbf{ML-KEM} (FIPS 203): Lattice-based key encapsulation,
    replacing ECDH for key exchange.
    \item \textbf{ML-DSA} (FIPS 204): Lattice-based digital signatures,
    replacing ECDSA/Ed25519 for signing.
    \item \textbf{SLH-DSA} (FIPS 205): Hash-based digital signatures,
    providing a conservative alternative based on minimal assumptions.
\end{enumerate}

% ============================================================================
